\chapter{Experimental Methodology and Data Analysis}
\label{ch:methodology}

Chapters~\ref{ch:data-models} and \ref{ch:implementation} defined the frozen
data, candidate models, deployed software and integration boundaries. This
chapter specifies how those artifacts are evaluated. Its purpose is to make
the unit of analysis, timing boundary, decision rule and execution status
explicit before numerical outcomes are reported in Chapter~\ref{ch:results}.
The distinction is important because the repository contains both completed
experiments and strengthened protocols prepared after limitations were found
in earlier work.

The five research questions require different evidence. Offline images can
compare predictive behaviour, but cannot establish sustained operation on an
edge device. A camera-free device replay can measure a processing path, but
cannot include acquisition latency. A stationary lighter trial can establish
that the integrated request--verdict path executed under one bounded
condition, but cannot estimate autonomous-flight reliability or industrial
transfer. The methodology therefore preserves each evidence category rather
than combining unlike observations into one performance claim.

\section{Study design and evidence hierarchy}
\label{sec:methodology:1}

The study uses a staged engineering-evaluation design. Each stage answers a
different part of the research questions and admits evidence only when its
configuration, raw records and acceptance decision are retained. The stages
are:

\begin{enumerate}
  \item \textbf{offline predictive evaluation}, using the frozen validation
  and test manifests to compare six detector configurations;
  \item \textbf{exact-device measurement}, using a fixed image replay to
  measure executable model--runtime--platform combinations;
  \item \textbf{stationary integrated evaluation}, using synchronized event
  identifiers and manually marked visible-flame intervals on the physical
  Sentinel and Scout;
  \item \textbf{application communication evaluation}, using correlated ROS~2
  requests and replies on the SwarmNet network;
  \item \textbf{mission simulation}, using two namespaced PX4/Gazebo vehicles
  and scenario-defined local-ENU truth; and
  \item \textbf{external proxy evaluation}, admitted only after licensing,
  annotation alignment and overlap checks pass.
\end{enumerate}

Evidence is ordered by traceability rather than by how favourable it appears.
Raw machine-readable records and immutable manifests are the primary
observations. Deterministic analysis outputs derived from those records are
calculated evidence. Tables and figures in Chapter~\ref{ch:results} transcribe
accepted derived evidence, while engineering interpretations are reserved for
Chapter~\ref{ch:discussion}. Screenshots and operator notes provide context,
but they do not replace a missing log, configuration or ground-truth interval.

Table~\ref{tab:methodology-status} records the execution boundary at this
chapter release. ``Executed'' means that the retained protocol produced an
accepted evidence bundle. ``Development only'' means that an executable path
was demonstrated but the planned repeated campaign was not completed.
``Pending'' means that the protocol exists but cannot contribute a result.

\begin{table}[H]
  \centering
  \caption{Experimental evidence classes and execution status at the Chapter 6
  release.}
  \label{tab:methodology-status}
  \scriptsize
  \begin{tabularx}{\textwidth}{@{}p{0.24\textwidth}p{0.18\textwidth}X@{}}
    \toprule
    Evidence class & Status & Admissible scope \\
    \midrule
    Six-model validation and frozen test & Executed & Predictive comparison on
    the curated internal data under the specified object-matching and operating
    point rules. \\
    Historical deployed-device replay & Executed & Camera-free processing
    measurements for the exact historical model, runtime, shape and device;
    not the strengthened six-model device matrix. \\
    Twelve model/platform device cells & Pending & The strengthened protocol is
    frozen, but no cell is admitted until conversion, fidelity and device gates
    pass. \\
    Concurrent-service endurance & Executed with limitations & A stationary
    30-minute complete-stack run combined both live perception paths, ROS~2,
    logging, Scout PX4 services and Sentinel Pixhawk/MAVProxy telemetry; the
    Scout latency and strict Jetson swap criteria were not met. \\
    Stationary integrated flame trials & Executed, bounded & Proof of the
    correlated physical perception workflow under marked laboratory exposure;
    no flight, localization or industrial-site inference. \\
    ROS~2 request--reply trials & Executed & Separate 300-event camera-free
    campaigns and a 1,000-event campaign under concurrent live inference load;
    all timing uses one Sentinel monotonic clock. \\
    Dual-vehicle mission scenarios & Development only & Executable nominal and
    selected fault paths with simulator-local truth; no formal reliability
    estimate. \\
    External proxy evaluation & Pending & No transfer result is admitted until
    archive, licensing, alignment and contamination gates close. \\
    \bottomrule
  \end{tabularx}
\end{table}

The mapping in Table~\ref{tab:rq-method-map} prevents an unanswered question
from receiving an inferred positive answer. It also identifies the experimental
unit used by each protocol. Images, source-family groups, application events,
physical exposures and simulation runs are not interchangeable independent
samples.

\begin{table}[H]
  \centering
  \caption{Research-question mapping to protocols, experimental units and
  final evidence disposition.}
  \label{tab:rq-method-map}
  \scriptsize
  \begin{tabularx}{\textwidth}{@{}p{0.07\textwidth}p{0.35\textwidth}p{0.25\textwidth}X@{}}
    \toprule
    RQ & Principal protocol & Unit of analysis & Status \\
    \midrule
    RQ1 & Frozen offline validation/test and exact-device screening & Object,
    image, source-family group; model--platform cell & Offline comparison
    executed; strengthened device matrix not executed. \\
    RQ2 & Full-service endurance with camera, inference, communication and
    logging active & Device run and time-indexed resource sample & Thirty-minute
    accepted complete-stack evidence exists; the run passed duration and
    process-survival gates but failed the Scout latency and strict swap gates. \\
    RQ3 & Common-stimulus stage ablation and synchronized stationary trials &
    Labelled frame sequence or marked physical exposure; correlated event &
    Integrated proof-of-function executed; comparative effect not evaluated. \\
    RQ4 & Preregistered two-vehicle simulation campaign & Retained run with a
    unique seed and attempt identifier & Development paths executed; formal
    310-run campaign not executed. \\
    RQ5 & Frozen-model inference on eligible external proxy datasets & Eligible
    image or independently defined video event & Admission work recorded;
    inference evaluation not executed. \\
    \bottomrule
  \end{tabularx}
\end{table}

\section{Offline predictive evaluation}
\label{sec:methodology:2}

The internal predictive comparison uses the curated fire-only manifest from
Chapter~\ref{ch:data-models}. Its validation split contains 700 images,
392 annotated fire boxes and 381 background images distributed across 307
reconstructed source-family groups. The one-time test split contains 733
images and 503 fire boxes. All six detector configurations use the same
manifests, label policy and prediction floor of 0.001. The validation data are
used to choose an operating threshold; the test data are evaluated once with
that frozen threshold and cannot be used for later tuning.

\subsection{Object matching and detection metrics}

Predictions are sorted by confidence and matched greedily to unmatched
ground-truth boxes in the same image. A pair is a true positive when its
intersection over union is at least 0.50. For predicted box \(B_p\) and
ground-truth box \(B_g\), intersection over union is
\begin{equation}
  \operatorname{IoU}(B_p,B_g)=
  \frac{|B_p\cap B_g|}{|B_p\cup B_g|}.
  \label{eq:method-iou}
\end{equation}
An unmatched prediction is a false positive and an unmatched ground-truth box
is a false negative. Precision, recall and F1 are then
\begin{align}
  P &= \frac{TP}{TP+FP}, &
  R &= \frac{TP}{TP+FN}, &
  F_1 &= \frac{2PR}{P+R}.
  \label{eq:method-prf}
\end{align}
The denominator is stated whenever a metric is reported. Object precision is
not described as an image-level false-alarm probability. For the background
subset, the reported background-image false-positive proportion is
\begin{equation}
  p_{\mathrm{BG,FP}}=
  \frac{\#\{\text{background images with at least one retained prediction}\}}
       {\#\{\text{eligible background images}\}}.
  \label{eq:method-bgfp}
\end{equation}
This proportion measures whether an image contains any retained false
detection. It is not a false-positive rate per hour or per video frame.

Average precision follows the COCO convention across IoU thresholds from 0.50
to 0.95 in steps of 0.05, with AP50 reported separately
\cite{Lin2014COCO}. These ranking metrics are threshold-integrated and are
therefore distinguished from the single operating-point counts used to compare
deployed decisions.

\subsection{Operating-point selection and scale analysis}

For each model, thresholds from 0.01 through 0.99 in increments of 0.01 are
evaluated on validation data. The selected threshold maximizes object-level
F1 at IoU 0.50. A tie is resolved first in favour of higher recall and then in
favour of the lower threshold. This deterministic rule prevents a later manual
choice based on the test set. The selected threshold and model hash are written
to the frozen evaluation record before the one-time test pass.

Recall is also stratified by ground-truth box area relative to image area:
less than 1\%, 1--5\%, 5--20\% and at least 20\%. The bins describe apparent
target scale in the curated images; they are not physical fire-size categories
because camera distance and calibration are unavailable. Laptop inference
timing uses batch size one after ten excluded warm-up iterations. It is a
software-screening measure and is not transferred to either edge platform.

Uncertainty intervals are calculated by a grouped nonparametric percentile
bootstrap. The resampling unit is the reconstructed source-family group, not
the individual image. Two thousand bootstrap samples use seed 20260729 and
resample the 307 validation groups with replacement. The 2.5th and 97.5th
percentiles form the descriptive 95\% interval. This follows the bootstrap
principle of resampling observed units \cite{Efron1979Bootstrap} while retaining
within-group dependence \cite{Harden2011ClusterBootstrap}. The intervals do
not correct annotation error, source-selection bias or imperfect group
reconstruction, and overlapping intervals are not treated as a formal
significance test.

\section{Exact-device screening and resource measurement}
\label{sec:methodology:3}

Two device-evidence programmes must remain separate. The first consists of
completed historical measurements of the deployed Sentinel and Scout
configurations. Those runs use their recorded runtimes, tensor interfaces,
model hashes and input shapes. Their camera-free replay begins with a fixed
\(480\times640\) black BGR image, applies the deployed preprocessing, executes
the runtime and decodes the output. Earlier baseline runs used five warm-up
and 50 measured iterations; later retained deployed measurements used their
own documented 100-iteration configuration. Each result is reported only with
its exact protocol.

The second programme is the strengthened cross-platform screen specified after
the six-model offline study. It defines 12 candidate cells: each of six models
on each of two edge platforms. A frozen, hashed sequence of 64
validation-calibration images is replayed in manifest order and cycled as
needed. Each admitted cell requires 20 warm-up iterations followed by 200
measured iterations at batch size one. Resource telemetry is sampled at 1~Hz.
The timed pipeline includes preprocessing, runtime execution and decoding;
camera acquisition, network transport and event coordination are excluded.

Each cell records source model hash, converted-artifact hash, runtime and
version, precision, input and output shapes, preprocessing, decoder, threshold,
device identity, power mode, cooling state and concurrent services. Conversion
success alone is insufficient. A fidelity check must establish that the
converted runtime behaves within the frozen acceptance rules on the calibration
set before latency or resource values are ranked. The final status vocabulary
is \texttt{DEPLOYED}, \texttt{TOO\_SLOW}, \texttt{OUT\_OF\_MEMORY},
\texttt{CONVERSION\_FAILED}, \texttt{RUNTIME\_UNSUPPORTED} or
\texttt{FIDELITY\_FAILED}. A failed cell remains part of the comparison rather
than disappearing from its denominator.

None of the 12 strengthened cells was admitted to the final evidence set. The
registered campaign was not executed before the evidence freeze.
Chapter~\ref{ch:results} may therefore report historical exact-deployment
measurements, but it cannot present them as the completed strengthened matrix.

\subsection{Timing boundaries}

Figure~\ref{fig:timing-boundaries} defines the timing quantities used across
the thesis. Runtime inference measures only the framework call. Device-pipeline
time includes preprocessing, runtime and decoding. The historical replay and
strengthened device screen are camera-free, so neither includes exposure,
driver buffering or image acquisition. Application round-trip time begins when
the Sentinel publishes an alert and ends when the matching verdict callback is
received on the Sentinel. Alert-to-decision time continues until the terminal
decision is recorded.

\begin{figure}[H]
  \centering
  \resizebox{\textwidth}{!}{\begin{tikzpicture}[
  font=\scriptsize,
  stage/.style={draw,rounded corners=2pt,align=center,text width=22mm,
    minimum height=12mm,fill=blue!4,inner sep=2mm},
  scout/.style={stage,fill=orange!9},
  event/.style={stage,fill=gray!10},
  flow/.style={-{Latex[length=2mm]},thick},
  measure/.style={|-|,thick},
  note/.style={align=center,font=\scriptsize}
]
  \node[stage] (capture) at (0,0) {camera exposure\\and acquisition};
  \node[stage] (pre) at (2.8,0) {preprocessing};
  \node[stage] (infer) at (5.6,0) {runtime\\inference};
  \node[stage] (decode) at (8.4,0) {decoding and\\thresholding};
  \node[event] (alert) at (11.2,0) {temporal gate and\\alert publication};
  \node[scout] (queue) at (14.0,0) {network, queue and\\Scout frame age};
  \node[scout] (verify) at (16.8,0) {Scout inference\\and verdict};
  \node[event] (decision) at (19.6,0) {Sentinel final\\decision};

  \draw[flow] (capture) -- (pre);
  \draw[flow] (pre) -- (infer);
  \draw[flow] (infer) -- (decode);
  \draw[flow] (decode) -- (alert);
  \draw[flow] (alert) -- (queue);
  \draw[flow] (queue) -- (verify);
  \draw[flow] (verify) -- (decision);

  \draw[measure] ([yshift=-12mm]infer.south west) --
    node[note,below=1mm]{runtime inference} ([yshift=-12mm]infer.south east);
  \draw[measure] ([yshift=-21mm]pre.south west) --
    node[note,below=1mm]{device pipeline: preprocessing + runtime + decoding}
    ([yshift=-21mm]decode.south east);
  \draw[measure] ([yshift=-30mm]alert.south west) --
    node[note,below=1mm]{application RTT: Sentinel publication to matching callback}
    ([yshift=-30mm]verify.south east);
  \draw[measure] ([yshift=-39mm]alert.south west) --
    node[note,below=1mm]{alert-to-decision duration on the Sentinel monotonic clock}
    ([yshift=-39mm]decision.south east);

  \node[note,anchor=north west,text width=49mm] at ([yshift=-49mm]capture.south west)
    {Camera exposure and acquisition are outside the camera-free backend replay.};
  \node[note,anchor=north east,text width=64mm] at ([yshift=-49mm]decision.south east)
    {Monotonic timestamps from different devices are never subtracted directly.};
\end{tikzpicture}
}
  \caption{Timing boundaries used by the evaluation. Brackets identify
  different measured quantities; they must not be compared as if they covered
  the same processing path.}
  \label{fig:timing-boundaries}
\end{figure}

For an ordered sample \(x_1,\ldots,x_n\), each latency or resource series is
summarized by sample size, arithmetic mean, median, minimum, maximum, range and
sample standard deviation
\begin{equation}
  s=\sqrt{\frac{1}{n-1}\sum_{i=1}^{n}(x_i-\bar{x})^2}.
  \label{eq:method-sd}
\end{equation}
Percentiles are added where the retained sample size supports them. Throughput
is computed from the complete measured interval rather than as the reciprocal
of a rounded mean latency. Warm-up observations are retained separately but
excluded from the reported steady measurement. No outlier is removed solely
because it is slow; exclusions require a recorded protocol violation or
invalid run state.

The two Sentinel engines require particular care. Historical Test~1 used a
416-pixel engine with tensors \( [1,3,416,416] \rightarrow [1,6,3549] \).
The final ROS workspace uses a distinct 640-pixel engine with tensors
\( [1,3,640,640] \rightarrow [1,6,8400] \). Latency or throughput from the
416-pixel artifact is never assigned to the deployed 640-pixel engine without
a new controlled measurement of the matching hash.

\section{Full-service endurance}
\label{sec:methodology:4}

RQ2 asks whether a selected configuration can operate sustainably with its
required concurrent services. A short backend loop cannot answer that
question. The registered endurance protocol therefore ran each device for at
least 30 minutes with the services required by its role active: camera input, deployed
preprocessing and inference, ROS~2 communication, asynchronous evidence
logging and the platform services needed for the test. The model hash, input
shape, runtime, power mode, cooling arrangement, ambient conditions and service
list must be frozen before the run.

Resource samples were requested at 1~Hz where the platform exposes the signal.
The retained fields include CPU load, GPU or accelerator load, resident and
available memory, temperature, power when a calibrated or platform-supported
reading is available, processing time, achieved throughput, camera-read
failures, inference errors, dropped log records and communication outcomes.
Unavailable measurements are marked unavailable; they are not estimated from
nominal hardware specifications.

Acceptance required completion of the full interval without process death,
unbounded memory growth, thermal shutdown, silent evidence loss or a safety
state violation. A time series is inspected for warm-up behaviour, sustained
drift and discrete failures. Temperature and processing time are plotted
against elapsed time, while resource summaries retain their units and sample
counts. A trend may be described, but a causal thermal-throttling claim requires
corresponding frequency or platform throttle evidence.

A later complete-stack run operated the live D435i/TensorRT Sentinel path, the
live Scout camera/TFLite path, ROS~2 transport, asynchronous logging, Scout PX4
services, Sentinel Pixhawk/MAVProxy telemetry and approximately 1~Hz
host-resource sampling for more than 30 minutes. The preregistered 360-alert
probe used a mismatched trial identifier and was correctly rejected by the
Scout. This setup failure is retained. A corrected 36-alert probe exercised
current-frame Scout inference during the final 180.245~s without treating the
synthetic alerts as camera-detection accuracy observations. Chapter
\ref{ch:results} therefore reports duration, service coexistence and resource
behaviour across the full interval, but request--verdict latency only for the
separate corrected three-minute probe.

\section{Temporal-filter and tiered-workflow evaluation}
\label{sec:methodology:5}

RQ3 requires a common-stimulus comparison of three decision stages:

\begin{enumerate}
  \item a single Sentinel frame at the frozen confidence threshold;
  \item the same Sentinel sequence with the three-consecutive-frame temporal
  filter; and
  \item the retained Sentinel alert followed by a correlated Scout verdict.
\end{enumerate}

A valid comparative unit must contain a complete ordered frame sequence or a
synchronized physical exposure with event-level ground truth. The same frames,
thresholds and eligibility rules are applied to all stages. The analysis then
counts ground-truth events retained, unsupported alerts removed and additional
decision delay introduced by filtering and verification. Event sensitivity is
the proportion of eligible hazard events that reach the relevant stage.
False-alert burden must use a declared exposure denominator, such as per
eligible negative event, per labelled negative sequence or per observed hour.
These denominators cannot be reconstructed from a handful of selected alert
records.

For the implemented filter, an observation qualifies when its highest retained
fire confidence reaches threshold \(\tau\). The run length resets after a
non-qualifying observation, and an alert is created when the run length reaches
three. Filter delay can be measured in processed observations. Conversion to
seconds requires the actual observation timestamps because the inference rate
is not assumed constant.

The accepted physical trials use stationary, propeller-free vehicles and a
handheld lighter positioned approximately 20--25~cm from cameras mounted at
approximately 19~cm for the Sentinel and 9.5~cm for the Scout. Ignition and
extinction intervals were marked separately from node startup. Correlated
event identifiers join the Sentinel alert, Scout verdict and Sentinel decision.
The analysis script verifies those joins and hashes the input logs before
writing a derived JSON result; raw records are not edited.

Only post-correction trials are eligible. Early Scout trials used an incorrect
image-decoding or output-decoding assumption and are retained as diagnostic
history. The later 0.05 Scout threshold was selected after diagnostic evidence,
so the physical trials using it are bounded engineering observations rather
than a preregistered accuracy estimate. The strong visible-flame exposure, the
small-flame exposure and six deliberately forced unlit requests answer
different questions and are reported separately. Forced unlit requests check
the correlated rejection path; they do not estimate an environmental
false-positive rate.

The older unsynchronized lighter recordings remain historical
proof-of-function only. Ignition was not synchronized to script or log start,
so script-start-to-detection time is invalid. Unlit intervals cannot be counted
as false negatives, and trial-level success rates cannot be inferred without
known exposure ground truth. Because the retained physical evidence is small,
heterogeneous and dependent, no confidence interval or significance test is
attached to it. The complete comparative tiered-workflow ablation was not executed.

\section{Communication and fault experiments}
\label{sec:methodology:6}

The completed communication protocol consists of three sequential,
camera-free ROS~2 trials with 100 request--response events each. The trials use
isolated domain identifiers 43, 44 and 45, Cyclone DDS bound to
\texttt{wlan0}, reliable volatile event topics, queue depth 100 and a 1-second
application timeout. The Sentinel publishes one request only after the previous
event has completed or timed out, which avoids uncontrolled queue growth but
also means that the observations form a serial sequence on one network setup.

Response completeness is
\begin{equation}
  C=\frac{N_{\mathrm{matching\ replies}}}{N_{\mathrm{published\ requests}}}.
  \label{eq:method-completeness}
\end{equation}
This is an application-level measure. It is not packet loss: DDS may retransmit,
and the experiment does not capture every link-layer packet. Application RTT
for event (i) is calculated entirely on the Sentinel monotonic clock,
\begin{equation}
  RTT_i=\frac{t^{S}_{\mathrm{callback},i}-
                    t^{S}_{\mathrm{publish},i}}{10^6}\;\mathrm{ms}.
  \label{eq:method-rtt}
\end{equation}
The superscript (S) indicates the Sentinel clock. Scout and Sentinel
monotonic epochs are never subtracted.

All 300 intended observations, including slow tails, remain in the summary.
The analysis reports the descriptive statistics defined in
Equation~\ref{eq:method-sd}, counts above 10, 20, 50 and 100~ms, and identifies
Tukey tail observations for inspection without deleting or winsorizing them.
No confidence interval or significance claim is made because the events are
serial, share one radio environment and do not represent independent network
deployments.

A separate loaded protocol publishes 1,000 sequential correlated events on an
isolated ROS~2 domain while both camera/inference nodes, asynchronous loggers,
resource samplers and an independent periodic tiered-workflow probe remain active.
The same Sentinel-clock RTT definition is used. Completeness receives a
two-sided Wilson interval because the outcome is a bounded delivery proportion;
latency remains descriptive because the serial events share one run and one
radio environment. The Sentinel Pixhawk is excluded from this loaded protocol,
so it does not become complete-stack evidence.

A separate 30-sample-per-device clock-offset probe uses midpoint estimates
based on a request and response timestamp. The minimum-RTT sample is retained
as the least-queued offset indication. This supports chronology checks only;
it does not justify retroactively subtracting unsynchronized cross-host wall
times for fine-grained latency.

Fault evidence is categorized by layer. The physical coordinator tests retain
duplicate, mismatched, stale and timed-out verdict behaviour. Simulation tests
exercise mission-level cases such as lost alert, Scout unavailable, delayed
verdict, sensor or inference error, low battery, invalid geofence and
duplicate/stale verdicts. A software unit test or mocked process handoff is not
reported as a physical-network or Gazebo fault result.

Four stationary hardware fault protocols use five unique event identifiers per
condition. Removing the Scout ROS camera interface tests fail-closed
\texttt{ERROR/no\_frame}; proving a live preflight and then stopping only the
bridge tests \texttt{ERROR/stale\_frame}. Removing the Scout verification
process tests the Sentinel's unchanged 10~s timeout. Finally, injecting two
\texttt{CONFIRMED} verdicts after each event has already finalized as
\texttt{TIMEOUT} tests idempotent rejection of late and duplicate outcomes.
These are deliberate interface or dispatch-boundary faults, not estimates of
natural fault frequency or wireless-delay distributions.

\section{Simulation and external-evaluation protocols}
\label{sec:methodology:7}

\subsection{Two-vehicle mission simulation}

The simulation unit is one retained run with a unique scenario, seed and
attempt identifier. Each run starts two namespaced PX4 v1.17 SITL instances and
the corresponding ROS~2 Jazzy mission processes. Vehicle identity, local-ENU
scenario coordinates, geofence, battery conditions, fault injection, expected
state path and terminal acceptance rule are fixed in the manifest. Simulator
truth is used only to assess the commanded mission path; it is not presented as
vision-based localization.

Development runs exercised nominal confirmed and rejected branches on
provisional vehicle dynamics, together with selected bounded faults. A later
accepted visual confirmed-hazard run used distinct X500 V2-derived Sentinel and
ModalAI/Starling-derived Scout proxies and retained synchronized mission events,
video and hashes. These records show that the software path can execute in that
development environment. The
formal campaign is a separate preregistered programme of 310 runs: 50 nominal
confirmed, 50 nominal rejected, 50 repeated-confounder runs and 20 repetitions
for each of eight fault scenarios. Seeds and run identifiers are unique and
deterministic. Automatic retry is prohibited; a rerun receives a new attempt
identifier so that failed attempts remain visible.

Formal acceptance requires the scenario preflight, expected state sequence,
terminal vehicle state, safety invariant checks, complete logs and artifact
hashes to pass. Selected runs also require visual captures, but video does not
replace telemetry. The accepted visual run closes that requirement for one
  confirmed-hazard pilot. A final qualitative review additionally accepted the
  continuous recording from the v17 run for presentation after checking role
  labels, continuity and framing; this disposition adds no mission trial to the
  denominator. In the final evidence set, the repeated campaign was not executed
because measured airframe parameters, the full fault set and the retained-run
matrix have not all passed. Consequently, development traces cannot be
converted into a mission success percentage.

\subsection{External proxy evaluation}

External evaluation freezes the internal models, preprocessing and operating
thresholds. A dataset is eligible only after five gates close: a compatible
licence is recorded; the acquired archive and extracted files are hashed; files
decode successfully or every failure is enumerated; labels or event intervals
can be mapped to the fire-only policy; and exact, perceptual and source-family
overlap with training, validation and test data is checked. The eligible
manifest is frozen before inference.

The FASDD publisher reports 25,097 UAV samples under CC BY-SA 4.0, but the
archive has not been acquired and admitted. Its sample count is therefore
dataset context rather than an evaluated denominator. The MmodalFire archive
and licence checks passed for 28 videos. Decoding produced 2,974,209 frames
against 2,974,228 container-reported frames, with 19 one-frame warnings. These
are readiness observations, not detector results. Condition and event alignment
is incomplete: the stricter timestamp consolidation validated eight of 28
ranges and left 20 for review. No external inference is admitted while those
gates remain open.

If admitted, video analysis defines independent events or condition intervals
before computing metrics; millions of neighbouring frames are not treated as
independent trials. Results are stratified by eligible dataset, acquisition
condition and apparent target scale. Any finding is described as transfer to a
specific UAV or industrial-like proxy dataset. It does not establish operation
at an industrial site.

\section{Statistics, reproducibility and exclusions}
\label{sec:methodology:8}

Every reported table must identify the artifact and configuration, sample
size, unit, threshold, timing boundary and evidence source. Raw observations
remain append-only. Derived files record their input paths and SHA-256 hashes,
analysis command or script version, configuration hash, creation time and
acceptance outcome. A correction creates a new derived artifact; it does not
silently edit the source log.

Data cleaning is rule-based. A record is excluded only when a predefined
eligibility condition fails or a documented protocol violation makes the
quantity uninterpretable. Examples include pre-correction Scout decoding,
missing formal identifiers, a mismatched model hash, an incomplete camera
preflight or a run outside its marked ground-truth interval. Runtime failures,
timeouts, conversion failures and out-of-memory outcomes are substantive
outcomes and remain in the experiment register. Missing values are reported as
missing rather than imputed from neighbouring samples.

Descriptive statistics are the default because several experiments are bounded
engineering studies rather than random samples from a defined population.
Confidence intervals are used only where the resampling unit and dependence
structure are defensible, as in the grouped bootstrap for internal predictive
metrics. Statistical significance is not claimed merely because two point
estimates differ or two intervals fail to overlap. For repeated simulation, a
binomial success proportion and interval may be reported only after the formal
campaign is admitted and the denominator includes every retained attempt.

The evidence cutoff for this final controlled revision is 30 September 2026,
after admission and author visual acceptance of the v17 mission delivery. Late evidence
entered only through coordinated revisions of
the affected methodology status, result table, research-question answer,
abstract and conclusion. This rule prevents a late number from being inserted
without updating the assumptions and limitations that give it meaning.

\reviewstatus{The offline six-model evaluation, historical deployed-device
measurements, loaded ROS~2 round-trip trial, 30-minute complete-stack endurance
run, accepted visual mission and four repeated hardware fault cases are
executed evidence. The strengthened 12-cell device screen, common-stimulus
ablation, formal 310-run simulation campaign and external transfer evaluation
were not executed. Chapter \ref{ch:results} preserves these boundaries and
does not infer a full positive answer for RQ1, RQ3, RQ4 or RQ5.}
